\phantomsection\addcontentsline{toc}{section}{KI-Verzeichnis}
\section*{KI-Verzeichnis}

Verzeichnis der verwendeten KI-Hilfsmittel

\begin{table}[htbp]
\centering
\small
\begin{tabularx}{\textwidth}{@{}c l Y@{}}
\toprule
\textbf{Nr.} & \textbf{KI-System} & \textbf{Einsatzzweck /\allowbreak{} Prompt} \\
\midrule
1 & Claude (Anthropic) & Erstellung der vollständigen Projektarbeit (Gliederung, Ausformulierung aller Kapitel, Abbildungen, Tabellen und Verzeichnisse) auf Basis der Vorlesungsskripte des Moduls ERP-Systeme und der vorgegebenen Beispielstruktur. Prompt: „Ich muss eine Projektarbeit für das Modul ERP schreiben, mein Thema ist: Analyse von Prozessineffizienzen durch Process Mining in ERP-Systemen: Eine Untersuchung der Datenstrukturen und Algorithmen am Beispiel des Order-to-Cash-Prozesses. Ich habe in dem Ordner alle Skripte gespeichert, schaue dir diese an und richte dich danach. Dazu habe ich auch eine Beispielstruktur der Projektarbeit reingelegt, halte dich zwingend daran. [\dots{}] Schreibe die komplette Projektarbeit und arbeite alles aus, was du musst.“ \\
2 & Claude (Anthropic) & Umstellung des Praxisteils von einem fiktiven Unternehmen auf einen synthetischen, mit PM4Py erzeugten und analysierten SAP-O2C-Event-Log einschließlich Durchführung der Simulation und der Process-Mining-Analysen (Discovery, Conformance Checking). Prompt: „Nimm kein reales Unternehmen, sondern öffentliche Event Logs (oft die beste Wahl). [\dots{}] kannst du statt eines Unternehmens einen öffentlichen Datensatz analysieren, z. B. die BPI Challenge Logs (etwa BPI 2019, ein realer Purchase-Order-Log eines Konzerns) oder einen synthetischen SAP-O2C-Log mit PM4Py. Das ist methodisch sauber, reproduzierbar und du brauchst keine Firmenfreigabe.“ \\
\bottomrule
\end{tabularx}
\end{table}

